% =====================================================
% part2_determinant.tex ― 第2部:行列式と階数の理論
% =====================================================

% =====================================================
\chapter{置換 ― 行列式を組み立てる部品}\label{chap:perm}
% =====================================================

第I部では行列式を $2$ 次・$3$ 次に限って導入した。
$4$ 次以上に進むには、サラスの方法のような「覚え方」ではなく、
すべての次数で通用する\textbf{定義}が必要になる。
その定義の部品となるのが\textbf{置換}である。
一見、行列と無関係な話が始まるように見えるが、
$2$ 次の式 $ad - bc$ の中にすでに置換が隠れていたことが、章末で明らかになる。

\section{置換とは}

\begin{teigi}[置換]
集合 $\{1, 2, \dots, n\}$ からそれ自身への全単射(一対一対応)
$\sigma$ を \textbf{$n$ 文字の置換}という。
$\sigma(1) = i_1, \dots, \sigma(n) = i_n$ であることを
\[
\sigma = \mat{1 & 2 & \cdots & n \\ i_1 & i_2 & \cdots & i_n}
\]
と書く。$n$ 文字の置換全体の集合を $S_n$ と書く。
\end{teigi}

要するに「$1$ から $n$ の並べ替え」である。
高校の場合の数で学んだとおり、並べ替えの総数は $n!$ 通りだから、
$S_n$ は $n!$ 個の要素をもつ。

\begin{rei}
$3$ 文字の置換は $3! = 6$ 個ある。たとえば
\[
\sigma = \mat{1 & 2 & 3 \\ 2 & 3 & 1}
\]
は「$1$ を $2$ に、$2$ を $3$ に、$3$ を $1$ にうつす」置換である。
\end{rei}

置換は写像なので\textbf{合成}ができる。$\sigma$ のあとに $\tau$ を行う合成を
$\tau\sigma$ と書く(行列の積と同じく右から先に作用する)。
また、並べ替えを元に戻す\textbf{逆置換} $\sigma^{-1}$ がつねに存在する。
合成しても置換、逆も置換、何もしない置換(恒等置換 $\varepsilon$)がある
――この構造は行列の積・逆行列・単位行列と瓜二つである。
このような「演算をもつ集合」の構造は\textbf{群}と呼ばれ、
大学数学の至るところに現れる。

\section{互換と符号}

\begin{teigi}[互換]
二つの文字 $i, j$ だけを入れ替え、他を動かさない置換を\textbf{互換}といい、
$(i \; j)$ と書く。
\end{teigi}

重要な事実は、\textbf{どんな置換も互換の積(繰り返し)で表せる}ことである。
これは日常の経験からも納得できる:どんな順番に並んだトランプの山も、
「$2$ 枚を選んで入れ替える」操作の繰り返しで好きな順に並べ替えられる。

\begin{rei}
$\sigma = \mat{1 & 2 & 3 \\ 2 & 3 & 1}$ は、
まず $1$ と $2$ を入れ替え、次に $1$ と $3$ を入れ替える
($\sigma = (1\;3)(1\;2)$、右から作用)ことで実現できる。
互換 $2$ 回である。
\end{rei}

互換による表し方は一通りではない。しかし、次の事実が成り立つ。

\begin{teiri}[偶奇の不変性]
一つの置換を互換の積で表すとき、表し方によらず、
\textbf{互換の個数の偶奇は一定}である。
\end{teiri}

そこで、偶数個の互換で書ける置換を\textbf{偶置換}、
奇数個で書ける置換を\textbf{奇置換}と呼び、置換の\textbf{符号}を
\[
\mathrm{sgn}(\sigma) =
\begin{cases}
+1 & (\sigma \text{ が偶置換}) \\
-1 & (\sigma \text{ が奇置換})
\end{cases}
\]
と定める。定義から $\mathrm{sgn}(\tau\sigma) = \mathrm{sgn}(\tau)\,\mathrm{sgn}(\sigma)$ が成り立つ。

\begin{itembox}[l]{深掘り:なぜ偶奇は表し方によらないのか}
この定理は「当たり前」ではない。証明の鍵は、$n$ 変数の\textbf{差積}
\[
\Delta = \prod_{1 \le i < j \le n} (x_j - x_i)
\]
を考えることである。置換 $\sigma$ で変数の添字を並べ替えると、
$\Delta$ は $+\Delta$ か $-\Delta$ のどちらかになる(因子の集合自体は
符号を除いて変わらないため)。そして\textbf{互換を一回行うごとに
$\Delta$ の符号はちょうど一回反転する}ことが直接計算で確かめられる。
もし同じ置換が偶数個の互換でも奇数個の互換でも書けたとすると、
$\Delta$ が $+\Delta$ かつ $-\Delta$ となり矛盾する。
「式の対称性を見張り役に使う」この論法は、5次方程式の解の公式が
存在しないことの証明(ガロア理論)にまでつながる、代数学の源流の一つである。
\end{itembox}

\begin{mondai}\label{q:perm}
$S_3$ の $6$ 個の置換をすべて書き出し、それぞれの符号を求めよ。
偶置換と奇置換は何個ずつあるか。
\end{mondai}

% =====================================================
\chapter{一般の行列式}\label{chap:gendet}
% =====================================================

\section{定義}

準備が整った。$n$ 次正方行列の行列式を一気に定義する。

\begin{teigi}[行列式]
$n$ 次正方行列 $A = (a_{ij})$ に対し、
\[
\det A = \sum_{\sigma \in S_n} \mathrm{sgn}(\sigma)\,
a_{1\sigma(1)}\, a_{2\sigma(2)} \cdots a_{n\sigma(n)}
\]
を $A$ の\textbf{行列式}という。
\end{teigi}

式の構造を言葉にすると:
「各行から $1$ 個ずつ、しかも\textbf{列が重ならないように}成分を選んで掛ける。
列の選び方(これが置換 $\sigma$)の全 $n!$ 通りについて、
符号 $\mathrm{sgn}(\sigma)$ を付けて足し合わせる」。

\begin{rei}[$2$ 次・$3$ 次との一致]
$n = 2$ のとき $S_2 = \{\varepsilon, (1\;2)\}$ で、符号は順に $+1, -1$。よって
\[
\det A = a_{11}a_{22} - a_{12}a_{21} = ad - bc
\]
となり、第I部の定義と一致する。$n = 3$ のときは $3! = 6$ 項あり、
偶置換 $3$ 個が $+$、奇置換 $3$ 個が $-$ を与えて、
サラスの方法の $6$ 項と正確に一致する(問 \ref{q:perm} の結果と
照らし合わせて確認してほしい)。
サラスの方法とは「$S_3$ の全要素を図形的に列挙する方法」だったのである。
$n = 4$ では $4! = 24$ 項となり、対角線の図では数え切れない。
だからサラスの方法は $3$ 次までしか使えなかったのだ。
\end{rei}

\section{行列式の基本性質}

定義から次の性質が導かれる。以下 $A$ は $n$ 次正方行列とする。

\begin{teiri}[基本性質]
\begin{enumerate}
\item \textbf{転置不変性}:$\det {}^t\!A = \det A$
(${}^t\!A$ は行と列を入れ替えた\textbf{転置行列})。
したがって「行」について成り立つ性質はすべて「列」についても成り立つ。
\item \textbf{多重線形性}:一つの行だけに注目すると、$\det$ はその行について
線形である。すなわち、第 $i$ 行が $\vect{u} + \vect{v}$ なら行列式は
第 $i$ 行を $\vect{u}$ にしたものと $\vect{v}$ にしたものの和に分かれ、
第 $i$ 行を $c$ 倍すれば行列式も $c$ 倍になる。
\item \textbf{交代性}:二つの行を入れ替えると行列式は $-1$ 倍になる。
特に、二つの行が等しい行列の行列式は $0$ である。
\end{enumerate}
\end{teiri}

多重線形性と交代性から、実用上もっとも重要な次の性質が出る。

\begin{teiri}[行基本変形と行列式]
\begin{enumerate}
\item ある行を $c$ 倍すると、行列式は $c$ 倍になる。
\item 二つの行を入れ替えると、行列式は $-1$ 倍になる。
\item \textbf{ある行に別の行の $c$ 倍を加えても、行列式は変わらない}。
\end{enumerate}
また、上三角行列(対角線より下がすべて $0$)の行列式は対角成分の積である。
\end{teiri}

性質 3 は、加えた後の行列式を多重線形性で二つに分解すると、
片方が「二つの行が比例する行列」の行列式となって $0$ になることから従う。

これで $4$ 次以上の行列式の\textbf{実用的な計算法}が手に入った:
\textbf{掃き出し法で上三角に変形し、対角成分を掛ける}
(途中の行交換の回数だけ符号を反転し、行のスカラー倍は記録しておく)。

\begin{rei}
\[
\det \mat{1 & 2 & 0 & 1 \\ 0 & 1 & 1 & 2 \\ 1 & 1 & 0 & 0 \\ 2 & 0 & 1 & 1}
\]
を計算する。第 $3$ 行から第 $1$ 行を引き、第 $4$ 行から第 $1$ 行の $2$ 倍を引く
(行列式は変わらない):
\[
\det \mat{1 & 2 & 0 & 1 \\ 0 & 1 & 1 & 2 \\ 0 & -1 & 0 & -1 \\ 0 & -4 & 1 & -1}
\]
第 $3$ 行に第 $2$ 行を加え、第 $4$ 行に第 $2$ 行の $4$ 倍を加える:
\[
\det \mat{1 & 2 & 0 & 1 \\ 0 & 1 & 1 & 2 \\ 0 & 0 & 1 & 1 \\ 0 & 0 & 5 & 7}
\]
第 $4$ 行から第 $3$ 行の $5$ 倍を引くと上三角になり、
\[
\det \mat{1 & 2 & 0 & 1 \\ 0 & 1 & 1 & 2 \\ 0 & 0 & 1 & 1 \\ 0 & 0 & 0 & 2}
= 1 \cdot 1 \cdot 1 \cdot 2 = 2.
\]
\end{rei}

\begin{itembox}[l]{深掘り:定義どおりの計算は宇宙の寿命でも終わらない}
定義の式は $n!$ 項の和である。$n = 20$ のとき
$20! \approx 2.4 \times 10^{18}$ 項。$1$ 秒に $10$ 億項処理できる計算機でも
約 $77$ 年かかる。$n = 25$ なら太陽系の年齢を超える。
一方、掃き出し法による計算はおよそ $n^3$ 回程度の演算で済み、
$n = 20$ なら $1$ 万回以下、一瞬である。
\textbf{「定義」と「計算法」は別物であり、よい理論はよいアルゴリズムを生む}
――この視点は数値計算・データサイエンスで一貫して重要になる。
\end{itembox}

\section{積の行列式と公理的特徴づけ}

\begin{teiri}[積の行列式]
$n$ 次正方行列 $A, B$ に対して
\[
\det(AB) = \det A \cdot \det B.
\]
\end{teiri}

第I部で「行列式は変換の面積(体積)拡大率」と述べたことを思い出すと、
この定理は「変換を続けて行えば拡大率は掛け算になる」という、
ごく自然な主張である。特に $A$ が正則なら $AA^{-1} = E$ の両辺の行列式をとって
\[
\det A^{-1} = \frac{1}{\det A}
\]
が従う。

\begin{itembox}[l]{深掘り:行列式は三つの性質で特徴づけられる}
実は次が成り立つ:行列を「$n$ 本の行ベクトルの組」とみなすとき、
\begin{enumerate}
\item 各行について線形(多重線形性)
\item 二つの行が等しければ値は $0$(交代性)
\item 単位行列 $E$ での値は $1$(正規化)
\end{enumerate}
の三つを満たす関数 $D$ は\textbf{行列式ただ一つ}である。
つまり行列式は「発見されるもの」ではなく、体積の性質(各辺について線形、
つぶれたら $0$、単位立方体は $1$)を式にしたら\textbf{一意に決まってしまうもの}なのである。

この特徴づけを使うと $\det(AB) = \det A \det B$ がエレガントに証明できる。
$A$ を固定し、$D(B) = \det(AB)$ という $B$ の関数を考えると、
$AB$ の各行は $B$ の対応する行だけで決まるから、$D$ は $B$ の行について
多重線形かつ交代的である。すると上の一意性から $D(B) = c \det B$
($c$ は定数)の形しかありえず、$B = E$ を代入して $c = \det A$ を得る。
成分計算を一切せずに証明が終わる――抽象的な特徴づけの威力である。
\end{itembox}

\begin{mondai}\label{q:det4}
(1) 上の例にならい、
$\det \mat{0 & 1 & 2 \\ 1 & 1 & 1 \\ 2 & 3 & 1}$ を行基本変形で計算せよ
(最初に行の入れ替えが必要なことに注意)。
(2) \textbf{ヴァンデルモンドの行列式}
$\det \mat{1 & a & a^2 \\ 1 & b & b^2 \\ 1 & c & c^2}$
を行基本変形により計算し、因数分解された形で表せ。
\end{mondai}

% =====================================================
\chapter{余因子展開と逆行列の公式}\label{chap:cofactor}
% =====================================================

\section{余因子展開}

行列式のもう一つの計算法として、「次数を一つ下げる」漸化式がある。

\begin{teigi}[余因子]
$n$ 次正方行列 $A$ から第 $i$ 行と第 $j$ 列を取り除いてできる
$n-1$ 次正方行列の行列式に符号を付けた
\[
A_{ij} = (-1)^{i+j} \det(\text{第 } i \text{ 行と第 } j \text{ 列を除いた行列})
\]
を $A$ の $(i,j)$ \textbf{余因子}という。
\end{teigi}

\begin{teiri}[余因子展開]
任意の第 $i$ 行に沿って
\[
\det A = a_{i1} A_{i1} + a_{i2} A_{i2} + \cdots + a_{in} A_{in}
\]
が成り立つ(任意の列に沿った展開も同様)。
\end{teiri}

符号 $(-1)^{i+j}$ は市松模様
\[
\mat{+ & - & + & \cdots \\ - & + & - & \cdots \\ + & - & + & \cdots \\ \vdots & \vdots & \vdots & \ddots}
\]
で覚えるとよい。$0$ の多い行や列に沿って展開すれば項数が減るので、
\textbf{行基本変形で $0$ を増やしてから展開する}のが実戦的な計算法である。

\section{余因子行列と逆行列の公式}

余因子展開には著しい副産物がある。
第 $i$ 行の成分を\textbf{別の}第 $k$ 行($k \neq i$)の余因子と組み合わせると
\[
a_{i1} A_{k1} + a_{i2} A_{k2} + \cdots + a_{in} A_{kn} = 0
\]
となる(これは「第 $k$ 行を第 $i$ 行で置き換えた、二つの行が等しい行列」の
行列式の展開とみなせるため)。展開公式とこの式をまとめると、
余因子を\textbf{転置して}並べた\textbf{余因子行列}
$\widetilde{A} = (A_{ji})$ に対して
\[
A\,\widetilde{A} = \widetilde{A}A = (\det A)\,E
\]
が成り立つ。よって次を得る。

\begin{teiri}[逆行列の公式]
$\det A \neq 0$ のとき、
\[
A^{-1} = \frac{1}{\det A}\,\widetilde{A}.
\]
\end{teiri}

$n = 2$ でこの公式を書き下すと、第I部で天下りに与えた
$2$ 次の逆行列公式がそのまま現れる。あの公式の正体は余因子行列だったのである。

\section{クラメルの公式}

\begin{teiri}[クラメルの公式]
$\det A \neq 0$ のとき、$A\vect{x} = \vect{b}$ の解は
\[
x_j = \frac{\det A_j(\vect{b})}{\det A}
\qquad (j = 1, \dots, n)
\]
で与えられる。ここで $A_j(\vect{b})$ は $A$ の第 $j$ 列を $\vect{b}$ で
置き換えた行列である。
\end{teiri}

\begin{itembox}[l]{深掘り:クラメルの公式は「使う」ものではなく「見る」もの}
クラメルの公式は美しいが、数値計算には不向きである。
$n$ 元の方程式を解くのに $n+1$ 個の $n$ 次行列式が必要で、
掃き出し法より桁違いに遅い。では無価値かといえば逆で、
この公式の真価は\textbf{解の性質が一目で読み取れる}ことにある。
たとえば「解 $x_j$ は係数と右辺の\textbf{多項式を $\det A$ で割ったもの}」
だと分かるから、係数を連続的に動かせば解も連続的に動く
($\det A \neq 0$ である限り)。経済学の比較静学
――パラメータが変化したとき均衡がどう動くか――では、
クラメルの公式で解を陽に書いて偏微分する、という使い方が標準的である。
\textbf{計算はアルゴリズムに、洞察は公式に}。役割分担を意識してほしい。
\end{itembox}

\begin{mondai}\label{q:cramer}
(1) $\det \mat{2 & 0 & 1 \\ 1 & 3 & -1 \\ 0 & 2 & 1}$ を第 $1$ 行に沿った
余因子展開で計算せよ。
(2) 第I部から登場している連立方程式
$2x + y = 5$, $x - 3y = -1$ をクラメルの公式で解け。
\end{mondai}

% =====================================================
\chapter{線形独立性と階数}\label{chap:rank}
% =====================================================

\section{線形独立と線形従属}

第I部の問 \ref{q:vec} で、あるベクトルを他のベクトルの線形結合で表す問題を
解いた。本章ではその逆向きの問い
――「与えられたベクトルたちに\textbf{無駄}はないか」――を考える。

\begin{teigi}[線形独立・線形従属]
ベクトル $\vect{a}_1, \dots, \vect{a}_k$ について、
\[
c_1 \vect{a}_1 + c_2 \vect{a}_2 + \cdots + c_k \vect{a}_k = \vect{0}
\]
を満たす実数の組が $c_1 = \cdots = c_k = 0$ \textbf{しかない}とき、
これらは\textbf{線形独立}であるという。
それ以外(すべては $0$ でない組が存在する)とき\textbf{線形従属}であるという。
\end{teigi}

線形従属とは、どれか一つのベクトルが残りの線形結合で書けてしまう
(=情報として無駄がある)ことと同値である。実際、たとえば $c_1 \neq 0$ なら
$\vect{a}_1 = -\dfrac{c_2}{c_1}\vect{a}_2 - \cdots - \dfrac{c_k}{c_1}\vect{a}_k$
と解ける。

幾何的には、$2$ 本のベクトルが線形従属とは同一直線上にあること、
$3$ 本が線形従属とは同一平面上にあることである。

判定は機械的にできる:$\vect{a}_1, \dots, \vect{a}_k$ を列に並べた行列を
$A$ とすれば、線形独立性の条件は
\textbf{同次方程式 $A\vect{c} = \vect{0}$ が自明な解 $\vect{c} = \vect{0}$
しかもたないこと}にほかならない。掃き出し法で判定できる。

\section{階数}

\begin{teigi}[階数]
行列 $A$ を行基本変形で簡約な階段形にしたとき、
零ベクトルでない行の本数(=主成分 $1$ の個数)を $A$ の\textbf{階数}といい、
$\mathrm{rank}\,A$ と書く。
\end{teigi}

簡約な階段形は一意に定まる(第I部で述べた)から、階数は矛盾なく定義される。
直観的には、階数とは\textbf{行列に含まれる「実質的な情報量」}である:
連立方程式でいえば実質的に独立な式の本数、
列ベクトルの言葉でいえば線形独立に取れる本数の最大値である。

\begin{teiri}[階数の性質]
$m \times n$ 行列 $A$ について、
\begin{enumerate}
\item $\mathrm{rank}\,A \le \min(m, n)$。
\item $A$ の行ベクトルから線形独立に選べる最大本数も、
列ベクトルから線形独立に選べる最大本数も、ともに $\mathrm{rank}\,A$ に等しい。
\end{enumerate}
\end{teiri}

\begin{itembox}[l]{深掘り:なぜ「行の独立本数」と「列の独立本数」が一致するのか}
行と列はまるで役割が違うのに、独立に取れる本数が必ず一致する
――線形代数で最初に出会う「不思議な定理」である。証明の骨子は二つの観察にある。
\begin{enumerate}
\item \textbf{行基本変形は、列ベクトルたちの間の線形関係を変えない}。
実際、変形を表す正則行列を $P$ とすると、$A$ の列 $\vect{a}_j$ の間の関係
$\sum_j c_j \vect{a}_j = \vect{0}$ は、両辺に $P$ を掛けても同値なまま保たれる。
\item 簡約な階段形では、主成分を含む列が線形独立で、
他の列はそれらの線形結合で書ける。よって列の独立本数 = 主成分の個数。
一方、零でない行は明らかに線形独立だから、行の独立本数も主成分の個数。
\end{enumerate}
両者が同じ「主成分の個数」を経由して一致する。
階段形という計算道具が、理論の証明の道具にもなっている点に注目してほしい。
\end{itembox}

正則性の同値条件(第I部の定理)も、階数の言葉で拡張できる。

\begin{teiri}[正則性の同値条件・拡張版]
$n$ 次正方行列 $A$ について、次はすべて同値である。
\begin{enumerate}
\item $A$ は正則である。
\item $\det A \neq 0$。
\item $\mathrm{rank}\,A = n$。
\item $A$ の列ベクトル($n$ 本)は線形独立である。
\item 同次方程式 $A\vect{x} = \vect{0}$ の解は $\vect{x} = \vect{0}$ のみ。
\item 任意の $\vect{b}$ に対し $A\vect{x} = \vect{b}$ はただ一つの解をもつ。
\end{enumerate}
\end{teiri}

この同値リストは今後も成長を続ける(第IV部で「$0$ は $A$ の固有値でない」が加わる)。
\textbf{一つの概念(正則性)を多くの角度から言い換えられること}こそが、
線形代数の見通しのよさの源泉である。

\begin{mondai}\label{q:indep}
(1) $\vect{a}_1 = \mat{1 \\ 2 \\ 3}$, $\vect{a}_2 = \mat{2 \\ 3 \\ 4}$,
$\vect{a}_3 = \mat{3 \\ 5 \\ 7}$ は線形独立か。従属ならば非自明な関係式を一つ示せ。
(2) $\mathrm{rank} \mat{1 & 2 & 3 \\ 2 & 4 & 6 \\ 1 & 1 & 1}$ を求めよ。
\end{mondai}

% =====================================================
\chapter{連立一次方程式の解の構造}\label{chap:structure}
% =====================================================

第I部の掃き出し法で、解は「一意・無数・なし」の三択だと知った。
本章はその完全な理論化であり、第2部のゴールである。

\section{同次方程式の解空間}

右辺が零ベクトルの方程式 $A\vect{x} = \vect{0}$ を\textbf{同次方程式}という。
同次方程式は必ず $\vect{x} = \vect{0}$ という解(自明解)をもつので、
問題は「それ以外の解があるか」である。

同次方程式の解の集合 $W$ には、際立った性質がある:
$\vect{u}, \vect{v} \in W$ ならば
\[
A(\vect{u} + \vect{v}) = A\vect{u} + A\vect{v} = \vect{0}, \qquad
A(c\vect{u}) = c\,A\vect{u} = \vect{0}
\]
だから、$\vect{u} + \vect{v} \in W$ かつ $c\vect{u} \in W$。すなわち
\textbf{解どうしの線形結合もまた解}である。
このように和とスカラー倍について閉じた集合を\textbf{部分空間}と呼ぶ
(第III部の主役の先取りである)。同次方程式の解集合を\textbf{解空間}という。

\begin{teiri}[解の自由度]
$A$ を $m \times n$ 行列とすると、同次方程式 $A\vect{x} = \vect{0}$ の解は
\[
n - \mathrm{rank}\,A
\]
個のパラメータを用いて、ちょうど一通りに表される。
この個数を解の\textbf{自由度}という。
\end{teiri}

理由は掃き出し法から直接読み取れる。簡約な階段形にしたとき、
主成分のある列に対応する未知数($\mathrm{rank}\,A$ 個)は他で決まり、
残りの未知数($n - \mathrm{rank}\,A$ 個)は自由に選べるパラメータになる。
\[
\text{(未知数の個数)} = \text{(縛られる個数)} + \text{(自由な個数)}
\]
という当たり前の勘定が、$n = \mathrm{rank}\,A + (\text{自由度})$ という
定理になっている。この式は第III部で\textbf{次元定理}として再登場する。

\section{非同次方程式:特殊解 + 同次解}

一般の $A\vect{x} = \vect{b}$($\vect{b} \neq \vect{0}$)の解集合は、
部分空間には\textbf{ならない}(たとえば $\vect{0}$ が解でない)。
しかし、同次方程式と美しい関係で結ばれている。

\begin{teiri}[解の構造定理]
$A\vect{x} = \vect{b}$ の一つの解(\textbf{特殊解})を $\vect{x}_0$ とすると、
すべての解は
\[
\vect{x} = \vect{x}_0 + \vect{w}
\qquad (\vect{w} \text{ は同次方程式 } A\vect{w} = \vect{0} \text{ の解})
\]
の形にちょうど一通りに書ける。
\end{teiri}

証明は一行である:$\vect{x}$ が解なら
$A(\vect{x} - \vect{x}_0) = \vect{b} - \vect{b} = \vect{0}$ だから
$\vect{w} = \vect{x} - \vect{x}_0$ は同次方程式の解。逆も明らか。

幾何的にいえば、同次方程式の解空間(原点を通る直線や平面)を、
特殊解 $\vect{x}_0$ の分だけ\textbf{平行移動}したものが非同次方程式の解集合である。
高校で学んだ「直線 $2x + y = 5$ は、原点を通る直線 $2x + y = 0$ を
平行移動したもの」という事実の、完全な一般化になっている。

\begin{rei}
\[
\begin{cases}
x_1 + x_2 + x_3 = 1 \\
x_1 + 2x_2 + 3x_3 = 2
\end{cases}
\]
を掃き出すと $\mathrm{rank}\,A = 2$、未知数 $3$ 個なので自由度は $1$。
$x_3 = t$ とおいて解くと
\[
\mat{x_1 \\ x_2 \\ x_3} = \mat{0 \\ 1 \\ 0} + t \mat{1 \\ -2 \\ 1}
\qquad (t \text{ は任意の実数}).
\]
第 $1$ 項が特殊解、第 $2$ 項が同次方程式の一般解である。
解集合は、方向ベクトル $\mat{1 \\ -2 \\ 1}$ の直線を
点 $\mat{0 \\ 1 \\ 0}$ を通るように平行移動したものである。
\end{rei}

\section{解の存在条件}

最後に「解なし」の場合を理論化する。

\begin{teiri}[解の存在条件]
$A\vect{x} = \vect{b}$ が解をもつための必要十分条件は
\[
\mathrm{rank}\,A = \mathrm{rank}\,(A \mid \vect{b})
\]
である($(A \mid \vect{b})$ は拡大係数行列)。
\end{teiri}

掃き出して「$0 = (\text{零でない数})$」の行が現れることは、
拡大係数行列側だけ階数が $1$ 増えることにほかならない、
というのがこの定理の中身である。また、列ベクトルの言葉では
「$\vect{b}$ が $A$ の列ベクトルたちの線形結合で書けること」と言い換えられる
――第 \ref{chap:rank} 章の冒頭で立てた問いに、これで完全な答えが付いた。

\begin{itembox}[l]{深掘り:この構造定理は微分方程式にもそっくり現れる}
「一般解 = 特殊解 + 同次方程式の一般解」という構造は、
連立一次方程式だけのものではない。たとえば線形微分方程式
\[
y'' - 3y' + 2y = e^{5t}
\]
の一般解は、右辺を $0$ にした同次方程式の一般解
$C_1 e^{t} + C_2 e^{2t}$ に、特殊解(たとえば $\frac{1}{12}e^{5t}$)を
足したものになる。線形漸化式でも同じである。
偶然ではない。微分も「和を保ち定数倍を保つ」操作、
すなわち\textbf{線形写像}だからである(関数の集合をベクトル空間とみなす
――この飛躍は第III部で正当化される)。
$A\vect{x} = \vect{b}$ で証明した一行の論法が、
微分方程式でも漸化式でも一字一句そのまま通用する。
\textbf{線形代数を学ぶとは、この「同じ構造」を見抜く目を手に入れることである}。
\end{itembox}

\begin{mondai}\label{q:struct}
(1) 連立方程式 $x + y = 1$, $2x + 2y = a$ が解をもつための
$a$ の条件を、階数の言葉で述べよ。
(2) 同次方程式
$x_1 + 2x_2 - x_3 + x_4 = 0$,
$2x_1 + 4x_2 - x_3 + 5x_4 = 0$
の解空間の自由度を求め、一般解をパラメータ表示せよ。
\end{mondai}
